\section{Forces}

Light and the forces are not added to matter. They are the other motions of the one tone, and the gauge structure of the Standard Model falls out of the dock, much of it measured. \cite{pollard2026vibetest}

The photon is the vector octet. The $8v$ part of the directions, excited as a long wave, is a massless field with two transverse polarizations and a gapless dispersion, charge-neutral so it streams free while charged knots stay bound. The same tone is matter when knotted and light when waving, one substance and two motions, and an earlier attempt to give the photon its own separate field was built and dropped because it broke that unity and was not needed. The vector octet is not only a photon. It also carries a non-abelian structure, the commutators $[\sigma_i, \sigma_j] = 2i\epsilon_{ijk}\sigma_k$ of the weak force measured directly in it.

The simplest motion of all is sound, a long slow wave of the tone spread across many docks, gapless and straight in its dispersion, the humblest emergent particle there is. The photon is the same kind of wave one step up, a chargeless phonon of the one field, which is why light and sound are kin, both the tone waving rather than knotted. The unity carries one caveat. The bare conserved mode, taken alone, spreads diffusively, wandering as the square root of time rather than racing along a sharp cone, so a truly relativistic massless particle, light or sound or a gravitational wave at the speed limit, appears to need a second conserved quantity, a momentum, beyond the single charge. The gauge photon is confirmed and its structure is right, but earning the sharp light-cone mode from the bare rule is the one named gap, taken up among the open questions.

The deeper structure is grand unification, forced rather than fitted. The Standard Model's symmetry does not fit inside the $\mathrm{SO}(8)$ of the bare directions, it is too small, and the smallest enlargement that holds it is reached by adding the ternary tone as one more axis, giving $\mathrm{SO}(10)$. So the tone is not an optional flavor. It is what the algebra demands to make room for matter. The sixteen-component spinor of $\mathrm{SO}(10)$ is exactly one generation, every quark and lepton with a right-handed neutrino in the slots the algebra supplies, the same $\mathrm{SO}(10)$ that grand unified theory reaches from the top down, here reached from the bottom up, from a tone on a dock. The octonion route mapped by Furey, Baez, and Dixon lands on the identical structure.

\begin{result}[Grand unification from the dock]
The Standard Model does not fit $\mathrm{SO}(8)$ but does fit $\mathrm{SO}(10)$, reached by adding the ternary tone, whose $16$-spinor is one full generation. The breaking through $\mathrm{SU}(5)$ gives $\sin^2\theta_W = 3/8$ at unification, with bottom-tau unification, $\mathrm{Tr}\,Y = 0$, and anomaly cancellation forcing the exact Standard-Model hypercharges. \cite{pollard2026vibetest}
\end{result}

Break $\mathrm{SO}(10)$ the way the geometry breaks it and real numbers appear. The weak mixing angle is one of the handful of pure numbers that nature fixes and experiments measure, the dial that sets how strongly the electromagnetic and weak forces blend. Most theories cannot produce such a number at all, only tell a story around it. This one delivers it by counting directions and nothing else, $\sin^2\theta_W = 3/8$ at the unification scale, exactly the value the leading grand-unified theories reach from the top down. That figure sits far above any experiment, and the standard quantum corrections carry it down to about $0.231$ at the energies we can reach, which is the number the experiments measure. Landing a measured constant of nature with nothing fitted is the kind of thing only a real structure does.

The anomalies that any consistent theory must cancel do so for only one assignment of charges, and it is exactly the Standard Model's, $1/6$ for the quark doublet and the rest, so the strange fractional charges of the quarks are forced rather than merely observed. Bottom-tau unification and a vanishing hypercharge trace follow as the same kind of consequence. The field that does the breaking is not an added scalar but a self-condensate, a bound pattern of the one tone, so even the Higgs is made of the same stuff as everything else.

The forces themselves behave correctly. Coupling a charged knot to the wave gives lattice quantum electrodynamics, with charge conserved to a part in a billion, the Gauss law holding, the Wilson loop measuring enclosed flux, and the Aharonov-Bohm phase exact. The non-abelian sector confines. A loop in an $\mathrm{SU}(2)$ field carries a real string tension at every coupling, and two opposite charges in one dimension feel a constant force that never lets them part, which is why a free quark is never seen. One number stays out of reach. The overall strength of the coupling, the fine-structure constant's worth, is left free by the bare law, reported open rather than fitted, though the three Standard-Model couplings, which miss each other on their own, very nearly meet at one high scale once the minimal supersymmetric content is added.

Taken together, electromagnetism, the weak force, the strong force, and gravity are four sectors of the one dock, so the zoo of particles and the zoo of forces are a single catalog, the ways one ternary tone can wave, knot, and bind on its twenty-four directions.
